\documentclass[a5paper,10pt]{article}

% https://github.com/InDevRus/filippov-solutions

% Нумерация страниц
\pagenumbering{gobble}

% Настройка полей
\usepackage{geometry}
\geometry{left=1cm}
\geometry{right=1cm}
\geometry{top=1cm}
\geometry{bottom=1.5cm}

% Вставка таблицы
\usepackage{array}
\usepackage{tabularx}

% Инструменты выравнивания формул
\usepackage{amsmath}
\usepackage{mathtools}

% Диакритика в математике
\usepackage{amsfonts}

% Вычисления размеров на ходу
\usepackage{calc}

% Удобные записи производных
\usepackage[italic=true]{derivative}

% Настройки сносок
\usepackage{enotez}

% Длинные знаки векторов
\usepackage{anyfontsize}
\usepackage[b]{esvect}

% Вставка картинок
\usepackage{graphicx}

% Вставка красной строки
\usepackage{indentfirst}
\setlength{\parindent}{1em}

% Условия в командах
\usepackage{xifthen}

% Луа-скрипты
\usepackage{luacode}

% Настройка команд отдельно в математическом и текстовом режимах
\usepackage{mathcommand}

% Для повторения LaTeX кода
\usepackage{multido}

% Конфигурация отступов формул
\delimitershortfall-1sp
\usepackage{mleftright}
\mleftright

% Растягивание объектов
\usepackage{relsize}

% Обтекание текстом
\usepackage{wrapfig}
\usepackage{subcaption}
\usepackage{floatflt}

% Поддержка ввода кириллицы
\usepackage[english, russian]{babel}

% Настройка корневой позиции для компилятора
\usepackage{import}
\providecommand*{\ProjectRootPath}{..}

\import{\ProjectRootPath/preambles/macros/}{theorems}

\import{\ProjectRootPath/preambles/fonts}{font_setup}

% Знак градуса
\usepackage{siunitx}

\import{\ProjectRootPath/preambles/macros/}{operators}
\import{\ProjectRootPath/preambles/macros/}{redefinitions}

\import{\ProjectRootPath/preambles/fonts}{letters_setup}
\import{\ProjectRootPath/preambles/fonts/adjustments}{custom_superscripts}

\import{\ProjectRootPath/preambles/custom}{formatting}
\import{\ProjectRootPath/preambles/custom}{frames}
\import{\ProjectRootPath/preambles/custom}{list_setup}

% TODO: перевести команды на современный стиль объявления

\begin{document}

$\tsys{
    Q\br{x; y} = y + \sqrt{1 + 2\pow{x}{2}}
    \\ P\br{x; y} = x + y + 1
}$;
$\tsys{
\sqrt{1 + 2\pow{x}{2}} = -y
\\ x = -y - 1
}$;
$\pow{y}{2} = 1 + 2\pow{y}{2} + 4y + 2$;
$\pow{y}{2} + 4y + 3 = 0$;
$\sub{y}{1} = -1$, $\sub{y}{2} = -3$, и оба являются корнями.

Особыми точками будут $\sub{A}{1}\br{0; -1}$ и $\sub{A}{2}\br{2; -3}$.

\begin{enumerate}
    \item $y = \phi - 1$. 
    $\der{\phi}{x}
    = \dfrac {\phi - 1 + \sqrt{1 - 2\pow{x}{2}}} {x + \phi}$; \\
    $\sub{\psi}{\phi}\br{x; \phi}
    = \sqrt{1 - 2\pow{x}{2}} - 1 
    = 1 + \dfrac {1} {2} \br{-2\pow{x}{2}} + o\br{\pow{x}{2}} - 1
    = O\br{\pow{x}{2}}
    = O\br{\pow{r}{2}}$,
    а поэтому, например, 
    $\dfrac {\sub{\psi}{\phi}\br{x; \phi}} {\pow{r}{1 + 0,5}} \to 0$
    при $r\br{x; \phi} \to 0$.
    
    Таким образом, уравнение с выделенными линейными частями имеет вид
    $\der{\phi}{x} = \linebreak = \dfrac {\phi} {x + \phi}$.
    $A = \begin{pmatrix} 1 & 1 \\ 0 & 1 \end{pmatrix}$.
    $\det\br{A - \lambda E} = \br{\lambda - 1}^2$.
    $\sub{\lambda}{1} = \sub{\lambda}{2} = 1$ и особая точка $\sub{A}{1}\br{0; -1}$
    является <<вырожденным узлом>>.
    
    $A - \sub{\lambda}{1} E
    = \begin{pmatrix} 0 & 1 \\ 0 & 0 \end{pmatrix} 
    \sim \begin{pmatrix} 0 & 1 \end{pmatrix}$,
    $\sub{v}{1} = \begin{pmatrix} 1 \\ 0 \end{pmatrix}$.
    Единственная сепаратриса $\sub{\phi}{1} = 0$ будет касательной для траекторий.
    
    \item $\tsys{x = \xi + 2 \\ y = \phi - 3}$.
    $\der{\phi}{\xi} = \dfrac {\phi + \sqrt{9 + 2\pow{\xi}{2} + 8\xi} - 3}
    {\xi + \phi}$.
    
    $\begin{aligned}
        \sqrt{9 + 2\pow{\xi}{2} + 8\xi} - 3
        & = 3 \sqrt{1 + \dfrac {2} {9}\pow{\xi}{2} + \dfrac {8} {9}\xi} - 3 \\
        & = 3\ \br{1 + \dfrac {1} {2} \br{\dfrac {2} {9} \pow{\xi}{2} + \dfrac {8} {9} \xi}
        + O\br{\br{\dfrac {2} {9} \pow{\xi}{2} + \dfrac {8} {9} \xi}^2}} - 3 \\
        & = \dfrac {4} {3} \xi + O\br{\pow{\xi}{2}} \\
        & = \dfrac {4} {3} \xi + O\br{\pow{r}{2}} \text{ при }r \to 0.\end{aligned}$
     
    Без нелинейностей получаем 
    $\der{\phi}{\xi} = \dfrac {\dfrac {4} {3} \xi + \phi} {\xi + \phi}$.
    $A = \begin{pmatrix} 1 & 1 \\ \frac {4} {3} & 1 \end{pmatrix}$. \\
    $\det\br{A - \lambda E} = \dfrac {1} {3} \br{3\pow{\lambda}{2} - 6\lambda - 1}$.
    $\sub{\lambda}{1,2} = \dfrac {3 \mp 2\sqrt{3}} {3}$, а особая точка $\sub{A}{2}\br{2; -3}$
    -- <<седло>>.    
    Для $\sub{\lambda}{1} = \dfrac {3 - 2\sqrt{3}} {3}$ имеем собственный вектор
    $\sub{v}{1} = \begin{pmatrix} \sqrt{3} \\ -2 \end{pmatrix}$,
    сепаратрису $\sub{\phi}{1} = -\dfrac {2} {\sqrt{3}} \xi$,
    а для $\sub{\lambda}{2} = \dfrac {3 + 2\sqrt{3}} {3}$
    имеем собственный вектор $\sub{v}{2} = \begin{pmatrix} \sqrt{3} \\ 2 \end{pmatrix}$,
    сепаратрису $\sub{\phi}{2} = \dfrac {2} {\sqrt{3}} \xi$.
\end{enumerate}

\end{document}
